% ===========================================================================
% Additions for main.tex — verified against the actual papers, August 2026.
% Every citation below was checked; none came from a chatbot's reference list.
% ===========================================================================

% --- BibTeX entries -------------------------------------------------------

@article{cui2025imp,
  title   = {Overcoming the Identity Mapping Problem in Self-Supervised
             Hyperspectral Anomaly Detection},
  author  = {Cui, Yongchuan and Zhang, Jinhe and Liu, Peng and
             Song, Weijing and Zeng, Yi},
  journal = {IEEE Transactions on Geoscience and Remote Sensing},
  year    = {2025},
  note    = {arXiv:2504.04115; code: github.com/yc-cui/Super-AD}
}

@article{huo2025msmnet,
  title   = {Multi-scale memory network with separation training for
             hyperspectral anomaly detection},
  author  = {Huo, Yu and Dong, Youqiang and Wang, Chenhao and
             Zhang, Min and Wang, Hai},
  journal = {Information Processing and Management},
  year    = {2025}
}

@article{haidar2025xaibs,
  title   = {Explainability-Driven Dimensionality Reduction for
             Hyperspectral Imaging},
  author  = {Haidar, Salma and Oramas, Jos\'e},
  journal = {arXiv preprint arXiv:2509.02340},
  year    = {2025}
}

@article{adbsnet2025,
  title   = {Reliable Task-Constrained Band Selection Method for
             Hyperspectral Anomaly Detection},
  journal = {Remote Sensing},
  volume  = {17}, number = {17}, pages = {3081}, year = {2025},
  doi     = {10.3390/rs17173081}
}

% --- Related Work paragraph -----------------------------------------------

Recent work has begun to characterise the failure modes of self-supervised
reconstruction detectors systematically. Cui \emph{et al.}~\cite{cui2025imp}
formalise the identity mapping problem, in which a sufficiently expressive
network eventually reconstructs anomalies as faithfully as background, and
organise existing remedies into perturbation, reconstruction, and
regularisation. Within that taxonomy MSNet is placed among methods that
modulate reconstruction using an error-derived mask. Notably, the authors
identify the choice of the coefficient balancing reconstruction against
regularisation as an unresolved difficulty shared across this family. The
present work addresses precisely that coefficient. Separately, MSMNet
\cite{huo2025msmnet} extends multiscale separation training with a memory
module and coarse-to-fine pseudo-label generation, which we treat as the
closest architectural successor to our backbone.

% --- Justification for the fixed-budget protocol ---------------------------
% Put this in Experimental Setup. It converts your leakage problem into a
% methodology contribution.

Reconstruction-based detectors do not converge monotonically in detection
performance: AUC typically peaks and then degrades as the network begins to
reconstruct anomalous pixels~\cite{cui2025imp}. We observe the same behaviour
on Coast, where AUC peaks near epoch~750 and declines thereafter
(Fig.~\ref{fig:convergence}). Selecting a checkpoint at the observed peak
would require access to ground-truth labels and therefore constitutes
model-selection leakage. We accordingly fix the training budget a priori at
$10 \times 150 = 1500$ epochs for every configuration and report the
final-iteration model, with mean $\pm$ standard deviation over five seeds.
All configurations share seeds, so per-seed differences are paired.

% --- Positioning the XAI contribution -------------------------------------
% Rewrite of the contribution claim. Note what it does NOT say.

Explanation-guided feature selection has been explored for supervised
hyperspectral classification, where post-hoc attributions from LRP, SHAP or
RISE are aggregated into per-band influence scores and used to construct
compact spectral subsets~\cite{haidar2025xaibs}. Attribution-like signals have
also been applied to band selection for anomaly detection, though existing
formulations rely on spectral attention and require a small number of labelled
anomalous samples~\cite{adbsnet2025}. Our use of attribution differs in
purpose rather than degree: rather than selecting inputs, we use the spatial
locality of Integrated Gradients attributions relative to the current soft
mask as a control signal for the separation coefficient. Because the residual
criterion can only increase this coefficient, whereas the attribution-locality
criterion can move it in either direction, the latter admits genuine
closed-loop control.

% --- Limitations (write this yourself, do not omit it) ---------------------

A reviewer will ask whether the alignment score is circular, since the soft
mask and the attribution both derive from the reconstruction error. The
distinction we rely on is that the mask records \emph{where} error occurs
while the attribution records \emph{what caused it}; these differ over the
receptive field, so that error arising from a pixel's own spectrum is
distinguishable from error arising from its spatial context. We validate this
empirically with a spatially permuted saliency control (Sec.~\ref{sec:controls})
and report the result whichever way it falls.
